\documentclass[a4paper,12pt]{article}
\usepackage[utf8]{inputenc}
\usepackage[italian]{babel}
\usepackage{listings}
\usepackage{xcolor}
\usepackage{float}
\usepackage{placeins}
\usepackage{url}
\usepackage{amssymb}
\setlength{\emergencystretch}{3em}

\lstset{
  language=C,
  backgroundcolor=\color{lightgray!20},
  frame=single,
  breaklines=true,
  numbers=left,
  numberstyle=\tiny,
  stepnumber=1,
  numbersep=5pt,
  captionpos=b,
  basicstyle=\ttfamily\small,
  keywordstyle=\color{blue}\bfseries,
  commentstyle=\color{green!60!black},
  stringstyle=\color{red},
  showstringspaces=false
}

\title{LabReSiD26 -- Hands-On 9}
\author{Davide Ferrara -- 518629}
\date{}

\begin{document}

\maketitle

% ---------------------------------------------------------------------------
\section{Esercizio 1: Analisi preliminare}
% ---------------------------------------------------------------------------

\subsection*{Domanda 1 -- Perché tre livelli di annidamento invece di uno solo}

% TODO

\subsection*{Domanda 2 -- Valore di \texttt{nlmsg\_len} dopo ogni chiamata}

% TODO

\subsection*{Domanda 3 -- Perché \texttt{chiudi\_attr()} va chiamato in ordine inverso}

% TODO

\subsection*{Domanda 4 -- Verifica che \texttt{rta\_len} di \texttt{IFLA\_LINKINFO} valga 52}

% TODO

% ---------------------------------------------------------------------------
\section{Esercizio 2: Implementazione}
% ---------------------------------------------------------------------------

\subsection*{Il codice \texttt{veth\_create.c}}

% TODO

\subsection*{Verifica}

% TODO

\end{document}
